\section{The New Warrior Monk}

\paragraph{The Role of the Knight}

In \emph{Gnosis}, Vol I, \textbf{Boris Mouravieff} describes the Knight of the Middle Ages:

\begin{quotex}
The elite man of the Middle Ages was the knight. Chivalry formed the nobility, the ruling class in that epoch where money did not yet hold the reins of public and private life. To be noble still meant to be disinterested. A nobleman of those times was characterized by his physical and muscular strength. He must be able to wear armour, and to handle heavy lances and swords. In his turn, deviations and abuses aside, the knight --- whose vigour and strength of arms made him master of those around him --- obeyed the orders of the Church. He was obliged to defend the weak and the oppressed, and to regulate public life, which was based on the work of the peasants and on craftsmanship.
\end{quotex}


\paragraph{Withered Emotional Life}


The nobility at that time was not intellectual, oftentimes not even educated; that was reserved to the clergy. This is in keeping with the nature of the warrior caste whose path is Bhakti yoga as we previously explained in Caste and Social Order\footnote{\url{https://gornahoor.net/?p=1698}}.

This sounds incredible to some, most likely because the average state of emotional development is quite restricted. What \textbf{Abbe Prevost} wrote some 300 years in his novel \emph{Manon Lescaut} applies all the more today:

\begin{quotex}
The ordinary run of men are susceptible to only five or six passions, within whose compass they live their whole lives, and to which all their emotions are reduced. Remove love and hate, pleasure and pain, hope and fear, and there is nothing left for them to feel. But people whose characters are more noble can be stirred in a thousand different ways; it is as if they had more than five senses, and could receive ideas and sensations that transcend the ordinary limits of human nature.
\end{quotex}


Perhaps this is a slight exaggeration, unless you consider that anger, anxiety, worry---all experienced by moderns---are forms of fear, sexual desire a form of pleasure, and so on. What is missing, however, are feelings of awe, wonder, and reverence. Or shame of one's nakedness. Or guilt about practices that used to be condemned by well-born and healthy minded men, yet are praised today.

The Knight had deep feelings of loyalty and fidelity to his superiors and his land, which he could never bring himself to broach. Who experiences that today, other than perhaps partisans? Yet they do not look out for the common good, but only for the good of their party.

\paragraph{Esoterism and Emotional Development}


There is a misconception that esoterism is fundamentally about the intellect, and that learning new facts, theories, philosophies, or systems represent stages along the path. However, in our time, we can see the problems of an over-developed intellect. Rather than leading us to a higher truth, it has in fact acted as a corrosive to an entire way of life in the West. It has cast doubt on all traditional wisdom, leading to atheism, materialism, scientism, and revolution as the touchstones of the intelligent and educated man of today.

While the material aspects of life may have improved, it has come at the cost of spiritual impoverishment. Since the feeling faculties have so deteriorated, the quest is to replace it with sex, drugs, and contrived thrills\footnote{\url{https://gornahoor.net/?p=1582}} through passive media such as movies, video games, and theme parks.

The intellect, on its own, cannot find any purpose to life. Like Schopenhauer, it sees the world as an absurdity, so it is crime to perpetuate it. Hence, Western men and women, especially the most intelligent and educated, are willing to contracept and abort themselves out of existence. They experience no passion about transmitting the inheritance\footnote{\url{https://gornahoor.net/?p=5454}} that has been bequeathed to them. While paying lip service to Darwinism, they promote a future of dysgenics. The Middle Age was saner, since they knew to subordinate the Intellect to Religion.

Hence, the goal of the esoteric path is to achieve balance among the physical, emotional, and intellectual parts of our inner life. Since our emotional life is the least developed of the three, it merits the most intense focus. In particular, negative emotions need to be overcome. It is the force of the right kind of passions that gives energy to the will. As we pointed out in Qualifications for Initiation\footnote{\url{https://gornahoor.net/?p=1161}}, the intellectual part of it is quite secondary. The real tests are tests of the will, not of intellectual knowledge.

The new elite will therefore need to have the emotional development of the Knight, yet be balanced. The elite will acquire the skills of administration, science, technology, medicine, philosophy, art, and so on, but he will be immune from the lures of the glamour of everyday life, particularly the love of money as the highest good, that entrap people today.

The esoteric path is just. It is the result of one's own efforts. Unlike worldly success, it cannot be faked.

\paragraph{Detachment and Prayer}


Apparently, there are two types of prayer, in marked contrast to each other. I will take Joel Osteen, a very effective preacher, as the exemplar of one of them. He says that you, too, can have fame, fortune, and a pretty wife, just like him, provided you let God act in your life. His message is so powerful and popular, that he has his own channel on satellite radio, so you can listen to him 24×7.

The other type is called \textbf{mental prayer}. The fruits of this path have been described this way by \textbf{Dom Vitalis Lehodey} in \emph{The Ways of Mental Prayer}:


\begin{quotex}
It detaches us from the world. It teaches us to make no account of the world's promises or threats, of its esteem or its contempt; for the world can neither make us happy nor virtuous; we are no better because it extols us to the clouds, nor any worse when it tramples us under foot. We are worth only just what we are worth in God's eyes. Mental prayer makes us dread the corruption of the world, the danger of its praises, the perfidiousness of its caresses, far more than its impotent fury. It makes us understand that God alone is to be considered; that no other's anger is to be feared, no other's esteem to be desired, that there is no other friendship on which the heart can securely rest.
\end{quotex}


You won't hear that on satellite radio, not even on the Catholic channel. So before you embark on an inner path, be certain that is what you want.


\flrightit{Posted on 2014-07-22 by Cologero}

\begin{center}* * *\end{center}

\begin{footnotesize}
\begin{sffamily}
\texttt{Constantine Aetos on 2014-07-27 at 15:31 said:}

Great article and I couldn't agree with you more. Unfortunately, the way things are in today's western society. The idea of a warrior monk is as dead as God is in the European psyche. Then again, only a few are destined to be warrior monks. I just hope there's enough of them that can fight against the soldiers of darkness.

\hfill

\texttt{William Neville on 2014-08-04 at 01:47 said:}

This is great food for thought. I published an article today that's related to some of the themes you explored here, although perhaps from more of a ``second estate'' perspective rather than a ``first estate'' one:
\url{http://www.radixjournal.com/journal/2014/8/3/reclaiming-the-way-of-the-sword}


\hfill

\texttt{William Neville on 2014-08-04 at 01:56 said:}

On an unrelated note, I recently rewatched the movie ``Equilibrium'', which made me think back to this article and its Mouravieff's emphasis that it is the emotional development of the Knight that differentiates him. The movie depicts a dystopian future in which the totalitarian government has taken the logic behind ``hate crime'' laws to its logical conclusion, banning and destroying art, poetry, and culture that provokes emotional reactions of any kind. The populace is prevented from feeling emotions via a drug that suppresses them and turns them into unfeeling, artificial stoics

Christian Bale plays an enforcer for this regime who comes to rediscover his feelings after executing his colleague (who he caught reading poetry that was supposed to have been disposed of). It's a bit of a B movie, but the political/spiritual messages are quite intriguing and sympathetic: \url{http://en.wikipedia.org/wiki/Equilibrium\_(film)}

\hfill

\texttt{Aleksandar on 2017-07-24 at 15:06 said:}

To me, there never was and probably never will be a more realizing, fulfilling role for a man than that of an idealized medieval knight. Warrior and monk, poet and philosopher, all in one, greater than life. A synergy of pure action and contemplation, a walking cathedral of God's Light.


\hfill

\end{sffamily}
\end{footnotesize}

